\documentclass[11pt]{article}
\usepackage[T1]{fontenc}
\usepackage{textcomp}
\usepackage[margin=0.75in]{geometry}
\usepackage{amsmath,amssymb,amsthm}
\usepackage{array,booktabs}
\usepackage{hyperref}
\setlength{\parindent}{0pt}
\newtheorem{theorem}{Theorem}
\theoremstyle{definition}
\newtheorem{definition}{Definition}
\title{Flow Proof Artifact: pythagoras-202121-witness}
\author{Generated by \texttt{flow doc proof}}
\date{Flow Proof Book}
\begin{document}
\maketitle
The 20-21-29 triangle satisfies the Pythagorean relation.\\[0.5em]
\textbf{Source.} euclid — Elements, Book I, Proposition 47\\[0.5em]
\bigskip
\noindent{\large \textbf{Derived fact 1.} \textit{29 squared equals 20 squared plus 21 squared} \hfill the Euclidean plane \cdot right triangle \cdot twenty twenty one twenty nine satisfies Pythagoras}\par
\smallskip\noindent\textit{Coordinate: the Euclidean plane \cdot right triangle \cdot twenty twenty one twenty nine satisfies Pythagoras}\par
\smallskip\noindent\textit{Source: euclid}\par
\smallskip\noindent\textit{Built on: seven twenty four twenty five satisfies Pythagoras, for right triangles in on the Euclidean plane, the square on the hypotenuse equals the sum of squares on the legs}\par
\medskip
\noindent\textbf{Goal.} 29 squared equals 20 squared plus 21 squared\par
\noindent$\displaystyle 29^{2} = 20^{2} + 21^{2}$\par
\medskip
\noindent
\begin{tabular}{@{} >{\raggedright\arraybackslash}p{0.06\textwidth} >{\raggedright\arraybackslash}p{0.47\textwidth} >{\raggedleft\arraybackslash}p{0.41\textwidth} @{}}
\textbf{\#} & \textbf{Proof} & \textbf{Mathematics} \\
\midrule
\textbf{1.} & We prove that twenty twenty one twenty nine satisfies Pythagoras for right triangles in the Euclidean plane. & \\[0.45em]
\textbf{2.} & We invoke the derived fact governing right triangles in the Euclidean plane: seven twenty four twenty five satisfies Pythagoras, for right triangles in on the Euclidean plane. & \\[0.45em]
\textbf{3.} & We invoke the derived fact governing right triangles in the Euclidean plane: the square on the hypotenuse equals the sum of squares on the legs. & \\[0.45em]
\textbf{4.} & From step 2 and step 3, this implies 29 times 29 equals 20 times 20 plus 21 times 21. Hence proven. & $\displaystyle 29^{2} = 20^{2} + 21^{2}$ \\[0.45em]
\bottomrule
\end{tabular}
\label{thm:Geometry-right triangle-twenty_twenty_one_twenty_nine_satisfies_pythagoras}
\par\medskip\hrule
\end{document}
